\documentclass[a4paper,17pt]{extarticle}
\usepackage{tikz,fontspec,makecell,graphicx,qrcode}
\usepackage[russian]{babel}
\usepackage[top=1cm,bottom=0cm,outer=1.8cm,inner=1.8cm]{geometry}
\parindent=0.6in \parskip=2.8mm

\setmainfont[
	Path = f/,
	BoldFont=lb.ttf,
	ItalicFont=li.ttf,
	BoldItalicFont=lbi.ttf
		]{l.ttf}
\setsansfont[
	Path = f/,
	BoldFont=lb.ttf,
	ItalicFont=li.ttf,
	BoldItalicFont=lbi.ttf
		]{l.ttf}

\begin{document} \thispagestyle{empty}

\begin{center} \begin{tabular}{ccc}
	\makecell[c]{\includegraphics[height=3.5cm]{spbtym-loc}} &
	\hspace{0.7cm} &
	\makecell[c]{\includegraphics[height=1.75cm]{funds}}
\end{tabular}

\begin{tabular}{l}
	\makecell[l]{{\large\bf С 22 по 28 марта мы приглашаем вас}\\
		{\large\bf принять участие в математической Регате}\\
		{\large\bf в рамках IX Санкт-Петербургского}\\
		{\large\bf Турнира юных математиков}}
\end{tabular}\end{center}

Регата — это командное соревнование, в котором командам будет предложено решить 3 серии исследовательских задач. На решение каждой серии у команд есть 2 дня. Решения оформляются в электронном виде и отправляются членам жюри для проверки. По результатам трёх туров будут распределены места команд, присуждены специальные премии и выданы призы. 

Особенности математической Регаты: \vspace{-7mm}
\begin{itemize} \itemsep=0mm
	\item[–] оригинальные исследовательские задачи;
	\item[–] являются идеальной подготовкой к формату Турнира;
	\item[–] решения проверяют авторы задач, что исключает ошибки в ходе проверки;
	\item[–] задачи связаны с интересными разделами математики и могут служить началом небольшой научной работы;
	\item[–] участие дистанционно, без привязки к конкретному месту и времени.
\end{itemize} \vspace{-9mm}

\begin{center}
	Собирайте команду и вступайте в бой!\\ {\bf Регистрация до 15 марта}
	\vspace{4mm}

	\qrcode[hyperlink,height=3.4cm]{https://forms.gle/WuuC2qttGcfETTUm8}
\end{center}

\tikz[remember picture,overlay]{
	\node[opacity=0.28,inner sep=0pt] at (current page.center)
	{\includegraphics[width=\paperwidth,height=\paperheight]{spbtym-page}};
%	\fill[white,opacity=0.7] (current page.center)
%		++(-0.5 * \paperwidth,-0.5 * \paperheight)
%		rectangle ++(\paperwidth,\paperheight);
}
\clearpage

\end{document}